% =====================================================================
%  Annexe C : Erreurs frequentes
% =====================================================================
\chapter{Erreurs fréquentes et solutions}
\label{ann:erreurs}

\chapitreresume{Cette annexe classe les erreurs que l'on rencontre le plus souvent
en apprenant Django. Pour chacune, elle donne le message typique, la cause probable
et la correction. Lisez toujours le message complet : Django indique presque
toujours le fichier et la ligne.}

\section{Lire un message d'erreur}

Un message Django contient généralement :
\begin{enumerate}
  \item le type d'exception, par exemple \code{TemplateDoesNotExist} ;
  \item la valeur ou le nom concerné ;
  \item la liste des chemins essayés, pour les gabarits ;
  \item la trace complète, avec les fichiers et les lignes.
\end{enumerate}

\begin{bonne}
Lisez la dernière ligne d'abord, puis remontez la trace. La dernière ligne donne
l'exception ; la première ligne de votre code dans la trace indique souvent
l'endroit à corriger. Ne copiez pas tout le message dans un moteur de recherche :
cherchez d'abord le type d'exception et le nom du fichier.
\end{bonne}

\section{Erreurs de projet et de configuration}

\begin{longtable}{L{5.2cm} L{4.5cm} L{4.7cm}}
\toprule
\textbf{Message} & \textbf{Cause} & \textbf{Correction} \\
\midrule
\endhead
\footnotesize
\code{ModuleNotFoundError} & Django absent de l'environnement virtuel
ou environnement inactif. & Activez \code{.venv}, puis
\code{pip install Django}. \\
\code{No installed app} & Le label \code{catalogue} est inconnu :
application absente de \code{INSTALLED\_APPS}. & Ajoutez
\code{"catalogue"} dans les réglages, ou \code{CatalogueConfig}
dans \code{apps.py}. \\
\code{ModuleNotFoundError} & Le paquet \code{catalogue} est
introuvable : dossier mal placé ou \code{\_\_init\_\_.py} absent. & Vérifiez l'arborescence et le fichier
\code{\_\_init\_\_.py}. \\
\code{ImproperlyConfigured} & Clé \code{SECRET\_KEY} absente ou vide. & Définissez
\code{SECRET\_KEY} dans \code{settings.py} ou l'environnement. \\
\code{DisallowedHost} & Domaine absent de \code{ALLOWED\_HOSTS}. & Ajoutez le
domaine, ou \code{localhost} et \code{127.0.0.1} en développement. \\
Port déjà utilisé & Un serveur tourne déjà. & Utilisez
\code{runserver 8001} ou arrêtez l'autre processus. \\
\code{ImportError} & Import circulaire ou nom \code{X} inexistant. &
Déplacez l'import dans une fonction, ou corrigez le nom. \\
\bottomrule
\end{longtable}

\section{Erreurs d'URL et de vue}

\begin{longtable}{L{5.2cm} L{4.5cm} L{4.7cm}}
\toprule
\textbf{Message} & \textbf{Cause} & \textbf{Correction} \\
\midrule
\endhead
\footnotesize
\code{NoReverseMatch} & Nom de route inconnu ou argument manquant. & Vérifiez
\code{app\_name}, le \code{name} de la route et les paramètres. \\
\code{Page not found (404)} & Aucune route ne correspond. & Vérifiez l'ordre des
\code{path} et le slash final. \\
\code{TypeError} & Une classe de vue est passée sans
\code{.as\_view()}. & Ajoutez \code{.as\_view()}. \\
\code{AttributeError} & \code{include} pointe vers un module sans
\code{urlpatterns}. & Créez \code{urls.py} avec
\code{urlpatterns}. \\
\code{404} sur un slug valide & Le convertisseur ne correspond pas. & Utilisez
\code{<slug:slug>} et vérifiez la valeur en base. \\
\code{405 Method Not Allowed} & Vue limitée à \code{GET} ou \code{POST}. &
Adaptez le décorateur ou envoyez la bonne méthode. \\
\code{CSRF verification failed} & Jeton absent ou cookie bloqué
(erreur 403). & Ajoutez
\code{\{\% csrf\_token \%\}} et vérifiez les cookies. \\
\bottomrule
\end{longtable}

\section{Erreurs de gabarit}

\begin{longtable}{L{5.2cm} L{4.5cm} L{4.7cm}}
\toprule
\textbf{Message} & \textbf{Cause} & \textbf{Correction} \\
\midrule
\endhead
\footnotesize
\code{TemplateDoesNotExist} & Chemin incorrect ou dossier absent de
\code{DIRS}. & Vérifiez \code{TEMPLATES} et le chemin exact. \\
\code{Invalid block tag} & Balise mal orthographiée ou non chargée. & Vérifiez
l'orthographe et ajoutez \code{\{\% load ... \%\}}. \\
\code{Invalid filter} & Filtre inconnu ou bibliothèque non chargée. & Chargez
\code{templatetags} avec \code{\{\% load biblio \%\}}. \\
\code{TemplateSyntaxError} & Balise non fermée, par exemple
\code{\{\% if \%\}} sans \code{\{\% endif \%\}}. & Fermez chaque bloc. \\
\code{VariableDoesNotExist} & Variable absente du contexte. & Vérifiez la clé
passée par la vue. \\
\code{NoReverseMatch} dans un gabarit & URL mal formée. & Corrigez le nom ou
l'argument de \code{\{\% url \%\}}. \\
Page vide à l'endroit attendu & Nom de variable incorrect. & Comparez le contexte
de la vue et le gabarit. \\
\bottomrule
\end{longtable}

\section{Erreurs de modèle et de base}

\begin{longtable}{L{5.2cm} L{4.5cm} L{4.7cm}}
\toprule
\textbf{Message} & \textbf{Cause} & \textbf{Correction} \\
\midrule
\endhead
\footnotesize
\code{no such table} & Migrations non appliquées. & Lancez \code{migrate}. \\
\code{no such column} & Modèle modifié sans migration. & Lancez
\code{makemigrations} puis \code{migrate}. \\
\code{NOT NULL constraint failed} & Champ obligatoire laissé vide. & Fournissez
une valeur ou autorisez \code{null}/\code{blank}. \\
\code{UNIQUE constraint failed} & Doublon sur un champ unique. & Vérifiez
l'existence avant \code{create}, ou gérez l'erreur. \\
\code{IntegrityError} & Clé étrangère invalide (contrainte
\code{FOREIGN KEY}). &
Vérifiez que l'objet lié existe. \\
\code{ProtectedError} & Suppression bloquée par \code{PROTECT}. & Supprimez
d'abord les objets liés, ou utilisez un autre comportement. \\
\code{DoesNotExist} & \code{get()} ne trouve rien. & Utilisez
\code{get\_object\_or\_404} dans une vue. \\
\code{MultipleObjectsReturned} & \code{get()} trouve plusieurs objets. & Filtrez
sur un champ unique. \\
\code{RelatedObjectDoesNotExist} & Profil un à un absent. & Créez le profil, ou
testez \code{hasattr}. \\
\bottomrule
\end{longtable}

\section{Erreurs de formulaire et de fichier}

\begin{longtable}{L{5.2cm} L{4.5cm} L{4.7cm}}
\toprule
\textbf{Message} & \textbf{Cause} & \textbf{Correction} \\
\midrule
\endhead
\footnotesize
Formulaire toujours invalide & Champ obligatoire vide ou erreur masquée. &
Affichez \code{form.errors} et vérifiez \code{cleaned\_data}. \\
\code{This field is required} & Champ non rempli. & Rendez-le facultatif avec
\code{required=False} ou corrigez la saisie. \\
Fichier non enregistré & \code{enctype} absent ou \code{request.FILES} oublié. &
Ajoutez \code{multipart/form-data} et le second argument. \\
\code{Upload a valid image} & Fichier corrompu ou Pillow absent. & Vérifiez
l'image et installez Pillow. \\
\code{KeyError: 'champ'} dans \code{clean()} & Champ déjà invalide. & Utilisez
\code{cleaned\_data.get}. \\
\code{ValueError} & Conversion d'une chaîne non numérique en entier. & Testez \code{isdigit()} ou utilisez un formulaire. \\
\bottomrule
\end{longtable}

\section{Erreurs de permission et d'authentification}

\begin{longtable}{L{5.2cm} L{4.5cm} L{4.7cm}}
\toprule
\textbf{Message} & \textbf{Cause} & \textbf{Correction} \\
\midrule
\endhead
\footnotesize
Redirection infinie vers login & \code{LOGIN\_URL} pointe vers une vue protégée. &
Vérifiez la route de connexion. \\
\code{403 Forbidden} après connexion & Permission absente. & Ajoutez la permission
au groupe ou à l'utilisateur. \\
\code{405} à la déconnexion & Lien au lieu d'un formulaire \code{POST}. & Utilisez
un formulaire avec \code{csrf\_token}. \\
\code{RelatedObjectDoesNotExist} & Profil membre absent pour
\code{user} (relation un à un). & Créez le profil
à l'inscription, ou testez \code{hasattr}. \\
Session perdue au redémarrage & \code{SECRET\_KEY} change ou base vidée. &
Stabilisez la clé et conservez la base. \\
Identifiants refusés & Nom d'utilisateur ou mot de passe invalide, ou
compte inactif. & Vérifiez \code{is\_active} et le mot de passe. \\
\bottomrule
\end{longtable}

\section{Erreurs de fichiers statiques}

\begin{longtable}{L{5.2cm} L{4.5cm} L{4.7cm}}
\toprule
\textbf{Message} & \textbf{Cause} & \textbf{Correction} \\
\midrule
\endhead
\footnotesize
CSS non chargé en développement & \code{STATICFILES\_DIRS} incorrect ou
\code{\{\% load static \%\}} absent. & Vérifiez le chemin et la balise. \\
\code{404} sur \code{/static/} & Fichier absent ou \code{DEBUG = False}. &
Vérifiez le fichier, ou lancez \code{--insecure} pour tester. \\
CSS non chargé en production & \code{collectstatic} non exécuté ou Nginx mal
configuré. & Lancez \code{collectstatic} et vérifiez l'\code{alias}. \\
Image média inaccessible & \code{MEDIA\_URL} ou \code{MEDIA\_ROOT} incorrect. &
Vérifiez les deux réglages et l'URL \code{static()}. \\
Ancien CSS affiché & Cache du navigateur. & Videz le cache ou utilisez un nom
hashé. \\
\bottomrule
\end{longtable}

\section{Erreurs de déploiement}

\begin{longtable}{L{5.2cm} L{4.5cm} L{4.7cm}}
\toprule
\textbf{Message} & \textbf{Cause} & \textbf{Correction} \\
\midrule
\endhead
\footnotesize
\code{500} en production, rien en local & \code{DEBUG = False} masque la trace. &
Consultez les journaux du serveur. \\
\code{Bad Request (400)} & \code{ALLOWED\_HOSTS} incomplet. & Ajoutez le domaine
exact, sans \code{https://}. \\
Fichiers statiques absents & \code{collectstatic} oublié. & Lancez la commande et
vérifiez \code{STATIC\_ROOT}. \\
\code{CSRF} échoue derrière un proxy & En-tête \code{X-Forwarded-Proto} non
transmis. & Configurez Nginx et \code{SECURE\_PROXY\_SSL\_HEADER}. \\
Connexion impossible à PostgreSQL & Identifiants ou service arrêté. & Vérifiez
\code{DATABASES} et l'état du service. \\
\code{Too many connections} & Fuite de connexions ou trafic élevé. & Utilisez le
pooling et surveillez la base. \\
\bottomrule
\end{longtable}

\begin{piege}
Ne corrigez jamais une erreur en désactivant une protection : \code{CSRF\_COOKIE\_SECURE},
\code{ALLOWED\_HOSTS}, les permissions ou \code{DEBUG}. Cherchez la cause réelle.
Une protection désactivée pour \og faire marcher \fg{} devient une faille en
production.
\end{piege}

\section{Méthode de débogage}

\begin{enumerate}
  \item Reproduire l'erreur de façon fiable.
  \item Lire le type d'exception et la dernière ligne de la trace.
  \item Vérifier la configuration : application installée, URL incluse, gabarit
        présent, migration appliquée.
  \item Ajouter un \code{print} ou un \code{breakpoint()} dans la vue pour voir
        les valeurs.
  \item Utiliser le shell Django pour tester la requête ORM isolément.
  \item Réduire le problème à un cas minimal.
  \item Écrire un test qui reproduit le bug, puis corriger.
\end{enumerate}

\begin{codebox}[python]{Inspecter une requête}
livres = Livre.objects.filter(titre__icontains="prince")
print(livres.query)
print(livres.explain())
\end{codebox}

\begin{note}
\code{str(queryset.query)} affiche le SQL généré. \code{explain()} montre le plan
d'exécution de la base. Ces deux outils sont précieux pour comprendre une requête
lente ou inattendue.
\end{note}
